problem:  
Let \(G\subset \mathrm{SO}(n+1)\) act by isometries on the round sphere \(S^{n}\) with cohomogeneity one.  Let the orbit types be \(G/H\) (principal) and \(G/K_\pm\) (singular) so that \(S^{n}/G \simeq [0,L]\) and the metric takes the standard \(G\)-invariant warped form  
\[
g = dt^{2} + \sum_{i=1}^{r} a_i(t)^{2}\, g_i,
\]
where each \(g_i\) is the \(G\)-invariant metric on an irreducible \(H\)-submodule \(V_i\subset T(G/H)\).

If \(\gamma\!:\!(0,L)\to (0,\infty)\) represents a \(G\)-invariant hypersurface (so that \(M_\gamma = G\cdot \gamma(t)\)\!), the Hsiang–Lawson reduction yields a functional for the area:
\[
\mathcal{A}(\gamma)=\int_0^L \Psi(t)\,\sqrt{1+\dot{\gamma}(t)^{2}}\, dt,\quad 
\Psi(t)=\prod_{i} a_i(t)^{m_i}.
\]
For the mean–curvature–penalized variant introduce
\[
\mathcal{F}_{\lambda}(\gamma)=\int_0^L \Psi(t)
\big( \sqrt{1+\dot{\gamma}(t)^{2}}  + \lambda\, |H_\gamma(t)|^{2}\sqrt{1+\dot{\gamma}(t)^{2}}\big)\, dt,
\qquad \lambda>0.
\]
Here \(H_\gamma(t)\) is the mean curvature expressed in terms of \(a_i,\dot{a}_i,\dot{\gamma},\ddot{\gamma}\).  
Near a singular orbit (say \(t=0\)) one has \(a_1(t)\!\sim\! t\), while the other \(a_i(0)\!\ne\!0\). Hence \(\Psi(t)\sim t^{m_1}\) and the Euler–Lagrange equation for \(\mathcal{F}_\lambda\) becomes a fourth‑order *singular* ODE with weight \(t^{m_1}\).

---

**Problem (Degenerate‑weight compactness for equivariant curvature energies):**  
For fixed small \(\lambda>0\), study the boundary‑value problem given by the Euler–Lagrange equation of \(\mathcal{F}_\lambda\) with smooth \(G\)-invariant hypersurface boundary conditions near \(t=0,L\).  

**Question:**  
Can the degeneracy of the weight \(\Psi(t)\) near \(t=0\) lead to *infinitely many* mutually non‑congruent \(G\)-invariant stationary hypersurfaces \(M_{\gamma_k}\subset S^{n}\), whose profile functions \(\gamma_k\) exhibit an increasing number of oscillations in \((0,\varepsilon)\) and cluster (in Hausdorff or varifold sense) onto the singular orbit \(G/K_-\)?  
If such accumulation cannot occur for any compact cohomogeneity‑one action on a sphere, determine the analytic obstruction—e.g., show that the weighted fourth‑order ODE admits a discrete, non‑accumulating set of regular fixed‑energy solutions.

---

Why is it a "good" problem:  
This problem now formulates explicitly the reduced functional including profile derivatives and the orbit‑space weight, and highlights a mathematically sharp issue:  
do curvature‑regularized variational ODEs with *singular weight coefficients* support an infinite sequence of symmetric stationary geometries, or does degeneracy fail to produce bubbling‑type behavior familiar from higher‑dimensional PDEs?

It lies at the intersection of:
- **Representation‑theoretic geometry:** the asymptotics of \(\Psi(t)\) encode how isotropy representations degenerate at singular orbits.  
- **Geometric analysis and singular ODE theory:** the fourth‑order weighted Euler–Lagrange equation tests the extension of compactness principles from regular to degenerate settings.  
- **Variational and stability theory:** the problem probes the landscape of equivariant critical points and the mechanism preventing or enabling their accumulation.

The statement is compact and conceptually clear—asking whether degeneracy near the singular orbit can destroy compactness—and any resolution would yield new insights.  
A positive answer would reveal a novel “equivariant oscillation cascade’’ phenomenon; a negative answer would isolate a fundamental compactness principle for curvature‑regularized symmetric energies.  Either outcome would connect geometry, analysis, and group actions in a way both deep and potentially field‑shaping, thus satisfying the criteria for a truly *good* research problem.